\begin{theorem}[Real L10 route sibling exhaustion]
\label{thm:real-l10-route-sibling-exhaustion}
Every accepted $\RealUp$ terminal seal read exposed to the L10 Real and
Completion route factors through exactly the displayed four sibling
certificates
$$
\begin{array}{c}
\autoref{ch:concrete-instances-dyadicratcore-namecert},\quad
\autoref{ch:concrete-instances-streamname-namecert},\quad
\autoref{ch:concrete-instances-regseqrat-namecert},\\
\autoref{ch:concrete-instances-real-namecert}.
\end{array}
$$
The route projects only the finite-window StreamName rows, the dyadic-radius
and tolerance ledger rows, the regular-sequence readback rows, and the local
$\RealUp$ terminal seal rows, together with their displayed $\BHist$,
$\hsame$, $\Cont$, $\Pkg$, and $\NameCert$ support.

No quotient stream, host real equality, selected limit, countable-choice
selector, ambient completeness coordinate, metric-completeness coordinate, or
standard bridge row is part of that accepted terminal seal read.
\end{theorem}
\begin{proof}
First read the selected transport package through
\autoref{thm:real-selected-transport-synchronized-package-from-stream-classifier}.
It supplies the finite-prefix selected full-shape package, the contextual
endpoint package, and the synchronized displayed e-one rows from a single
stream classifier source.

Next use
\autoref{thm:real-selected-transport-fit-synchronized-from-stream-classifier}.
The fit theorem identifies the same selected data as a synchronized
transported $\RealUp$ observation certificate; the selected index, prefix
surplus, four tails, context rows, continuation rows, and derived prefix row
remain explicit finite data.

The L10 sibling route children then identify where those rows can be consumed.
\autoref{thm:real-l10-open-phase-sibling-route} places every accepted
open-phase seal read over the StreamName, RegSeqRat, DyadicRatCore, and local
Real surfaces. \autoref{thm:realup-open-phase-four-face-pullback} and
\autoref{thm:real-open-phase-four-object-terminal-route} make this surface
deterministic and exclude deletion or replacement of any displayed face.
\autoref{thm:real-l10-four-face-exit-criterion} carries the same four-face
route to the current terminal exit read.

Finally project the four rows named in the statement. The StreamName row is the
finite observation window; the DyadicRatCore row is the dyadic endpoint and
radius ledger; the RegSeqRat row is the finite regular-sequence request and
readback surface; the local $\RealUp$ row supplies the terminal seal face after
those three rows are read. Since the cited route theorems expose no constructor
for quotient streams, host equality, selected limits, countable choice, ambient
completeness, metric completeness, or a standard bridge row, the displayed four
sibling certificates exhaust the accepted terminal seal route.
\end{proof}

\begin{theorem}[Real L10 terminal four-face support minimality pullback]
\label{thm:real-l10-terminal-four-face-support-minimality-pullback}
Let $\mathcal T$ be an accepted $\RealUp$ terminal seal read for the current
Phase Real and Completion exit route. If a Real-facing terminal presentation
has the same carrier verdict and the same classifier verdict as $\mathcal T$,
then that presentation pulls back to exactly the four sibling certificates
$$
\DyadicRatCoreUp,\quad \StreamNameUp,\quad \RegSeqRatUp,\quad \RealUp
$$
and to the displayed support rows
$$
\BHist,\quad \hsame,\quad \Cont,\quad \Pkg,\quad \NameCert.
$$
The pullback is minimal at the terminal seal: deleting any one of those four
sibling certificates, or deleting any one of the displayed support rows,
prevents the presentation from carrying the same Real-facing carrier and
classifier verdicts.

No quotient stream equality, selected limit, countable-choice selector, ambient
$\RealUp$ completeness coordinate, metric-completeness coordinate, aggregate
formalstatus row, or bridge row belongs to this terminal support surface.
\end{theorem}
\begin{proof}
First read the accepted terminal seal through
\autoref{thm:real-l10-route-sibling-exhaustion}. It places every accepted
$\RealUp$ terminal seal read on the four sibling certificates
$\DyadicRatCoreUp$, $\StreamNameUp$, $\RegSeqRatUp$, and $\RealUp$, carried by
the displayed $\BHist$, $\hsame$, $\Cont$, $\Pkg$, and $\NameCert$ rows.

Next apply
\autoref{thm:streamname-terminal-four-face-support-minimality} to the
StreamName-facing terminal section. The same carrier and classifier verdicts
cannot be presented after deleting the finite StreamName window, the dyadic
ledger face, the RegSeqRat readback face, the local Real seal face, or any one
of the listed support rows.

Finally compare the RegSeqRat terminal boundary by
\autoref{thm:regseqrat-real-terminal-four-face-exactness}. That exactness
result reads the local $\RealUp$ terminal seal precisely from the compatible
four-face packet over the shared support rows, so any Real-facing presentation
with the same carrier and classifier verdict must pull back to that packet and
to no smaller packet.

The exclusion list is inherited from the three cited route descriptions. Their
terminal packets contain no quotient stream equality, selected limit,
countable-choice selector, ambient $\RealUp$ completeness coordinate,
metric-completeness coordinate, aggregate formalstatus row, or bridge row, so
none of those rows can replace a missing sibling certificate or displayed
support row at the terminal seal.
\end{proof}

\begin{theorem}[Real L10 terminal route support induction]
\label{thm:real-l10-terminal-route-support-induction}
Let
$$
\mathcal T_0,\mathcal T_1,\ldots,\mathcal T_n
$$
be a finite replay chain of accepted Real-facing terminal presentations for
the current Phase Real and Completion exit route. Suppose each successive
replay step preserves the $\DyadicRatCoreUp$ endpoint and radius ledger, the
$\StreamNameUp$ terminal-window row, the $\RegSeqRatUp$ finite-request and
readback row, the local $\RealUp$ terminal seal row, and the displayed
$\BHist$, $\hsame$, $\Cont$, $\Pkg$, and $\NameCert$ support rows. Then the
terminal read of every $\mathcal T_i$ lies on the same four-face support
surface
$$
\DyadicRatCoreUp,\quad \StreamNameUp,\quad \RegSeqRatUp,\quad \RealUp
$$
with those displayed support rows. If a replay step fails to preserve the
terminal read, the failure is witnessed by one of the four faces or by one of
the displayed support rows.

The induction introduces no quotient stream, host real equality, selected
limit, countable-choice selector, ambient completeness coordinate,
metric-completeness coordinate, aggregate formalstatus row, or bridge row.
\end{theorem}
\begin{proof}
For the base presentation $\mathcal T_0$, apply
\autoref{thm:real-l10-terminal-four-face-support-minimality-pullback}. The
terminal read pulls back to the four sibling certificates and to the displayed
$\BHist$, $\hsame$, $\Cont$, $\Pkg$, and $\NameCert$ rows, and no smaller or
off-route packet carries the same Real-facing verdicts.

For the successor step, assume $\mathcal T_i$ has already been read on that
support surface. The replay hypotheses identify the next presentation
$\mathcal T_{i+1}$ on the same dyadic endpoint and radius ledger, the same
StreamName terminal window, the same RegSeqRat finite request and readback
row, the same local Real seal row, and the same displayed support rows.

Apply
\autoref{thm:streamname-terminal-four-face-support-replay-induction} to replay
the terminal-window support through the ordered four faces. Then use
\autoref{thm:regseqrat-real-seal-finite-window-witness-induction} to carry the
finite request and readback witness across the RegSeqRat-to-Real seal
boundary. Finally,
\autoref{thm:real-l10-route-sibling-exhaustion} says that every accepted
Real-facing terminal read on this route is exhausted by exactly those sibling
faces and rows.

Thus the successor presentation remains on the same support surface. If the
successor replay fails, one of the hypotheses just used is absent: the dyadic
ledger, StreamName window, RegSeqRat request/readback face, local Real seal,
or one of the displayed $\BHist$, $\hsame$, $\Cont$, $\Pkg$, and $\NameCert$
rows has changed or disappeared.

Finite induction over the replay chain now gives the claim for all
$\mathcal T_i$. The exclusion list is inherited at each step from the
minimality pullback and route exhaustion theorems, so no quotient stream, host
real equality, selected limit, choice selector, completeness coordinate,
formalstatus row, or bridge row can enter the terminal support surface.
\end{proof}

\begin{theorem}[Real L10 terminal support cover induction]
\label{thm:real-l10-terminal-support-cover-induction}
Let $\mathcal{U}$ be a finite family of accepted Real-facing terminal
presentations for the current Phase Real and Completion exit route. Suppose
each replay step inside each member of $\mathcal{U}$ preserves the displayed
$\DyadicRatCoreUp$ endpoint and radius ledger, the $\StreamNameUp$
terminal-window row, the $\RegSeqRatUp$ finite-request and readback row, the
local $\RealUp$ terminal seal row, and the shared $\BHist$, $\hsame$, $\Cont$,
$\Pkg$, and $\NameCert$ support rows. Suppose also that every finite subcover
of $\mathcal{U}$ satisfies the objectwise Phase Real and Completion threshold,
and that overlaps agree on those four faces and support rows.

Then the whole replay-cover satisfies the same objectwise Phase Real and
Completion threshold. If a member or an overlap fails to contribute to that
threshold, the failure localizes to one displayed face
$$
\DyadicRatCoreUp,\quad \StreamNameUp,\quad \RegSeqRatUp,\quad \RealUp
$$
or to one displayed support row among $\BHist$, $\hsame$, $\Cont$, $\Pkg$, and
$\NameCert$.

The cover induction asserts no quotient stream equality, selected limit,
countable-choice selector, metric-completeness datum, aggregate formalstatus
row, bridge row, or ambient $\RealUp$ completeness claim.
\end{theorem}
\begin{proof}
First run the terminal support induction inside each cover member. Apply
\autoref{thm:real-l10-terminal-route-support-induction} to each finite replay
chain. The result keeps every accepted Real-facing terminal presentation on
the same displayed $\DyadicRatCoreUp$, $\StreamNameUp$, $\RegSeqRatUp$, and
local $\RealUp$ support surface, with the shared $\BHist$, $\hsame$, $\Cont$,
$\Pkg$, and $\NameCert$ rows.

Next compare the StreamName and RegSeqRat cover faces. Apply
\autoref{thm:streamname-terminal-four-face-support-replay-induction} to keep
the terminal-window rows stable under replay. Apply
\autoref{thm:regseqrat-four-face-finite-cover-induction} to transport the
finite-request, readback, and regularity rows through the finite four-face
cover. The overlap hypotheses are exactly the rows consumed by these two
induction principles.

Now glue the local threshold packets. Apply
\autoref{thm:real-l10-finite-cover-exit-threshold-gluing} to the finite
family whose members already satisfy the objectwise Phase Real and Completion
threshold on the displayed faces. The overlap agreement supplies the common
dyadic ledger, StreamName window, RegSeqRat readback, local Real seal,
$\BHist$, $\hsame$, $\Cont$, $\Pkg$, and $\NameCert$ support needed by the
gluing theorem.

Finally descend from the glued finite packet to the whole replay-cover. Apply
\autoref{thm:real-l10-finite-cover-exit-threshold-descent}. Any failed
descent hypothesis is one of the same visible faces or support rows, so the
failure localizes there and not to an off-route completion coordinate.

The exclusion clauses in the support induction, the StreamName and RegSeqRat
cover inductions, and the gluing and descent theorems all refuse quotient
stream equality, selected limits, countable choice, metric completeness,
aggregate formalstatus rows, bridge rows, and ambient $\RealUp$ completeness.
Thus the composed cover result remains a finite four-face terminal support
claim and carries no route-level strengthening.
\end{proof}

\begin{theorem}[Real L10 route-cover phase-exit exactness]
\label{thm:real-l10-route-cover-phase-exit-exactness}
Let $\mathcal{U}$ be a finite replay-cover of accepted Real-facing terminal
presentations for the current Phase Real and Completion exit route. Suppose
$\mathcal{U}$ satisfies the terminal support cover induction of
\autoref{thm:real-l10-terminal-support-cover-induction}, and suppose the four
roadmap faces of every cover member and overlap meet the local scoped and
audit records required by
\autoref{thm:phase-real-completion-exit-four-object-threshold-projection}.

Then $\mathcal{U}$ supplies the public Phase Real and Completion exit read
exactly through
$$
\DyadicRatCoreUp,\quad \StreamNameUp,\quad \RegSeqRatUp,\quad \RealUp
$$
and through the same-route $\RealCompletionExactBoundaryUp$ witness. The read
uses the displayed $\BHist$, $\hsame$, $\Cont$, $\Pkg$, and $\NameCert$
support rows from the Real L10 route, and it admits no quotient stream,
selected limit, countable-choice selector, host-completeness row,
metric-completeness coordinate, aggregate formalstatus row, or bridge token.
\end{theorem}
\begin{proof}
First apply \autoref{thm:real-l10-terminal-support-cover-induction}. The
finite replay-cover is reduced to the displayed terminal support surface:
$\DyadicRatCoreUp$, $\StreamNameUp$, $\RegSeqRatUp$, local $\RealUp$, and the
shared $\BHist$, $\hsame$, $\Cont$, $\Pkg$, and $\NameCert$ rows. Any failed
cover contribution is already localized to one of those rows.

Next apply \autoref{thm:real-l10-route-sibling-exhaustion} and
\autoref{thm:real-l10-terminal-four-face-support-minimality-pullback}. These
results keep the Real-facing terminal read on the same L10 route and forbid a
smaller packet, off-route packet, quotient stream, selected limit, choice
selector, host-completeness row, metric-completeness coordinate, aggregate
formalstatus row, or bridge token from carrying the same verdicts.

Now read the phase-exit side. The hypotheses give the local scoped closure
records and audit-clean verification records for the four roadmap faces, so
\autoref{thm:phase-real-completion-exit-four-object-threshold-projection}
projects the accepted packet to exactly
$$
\DyadicRatCoreUp,\quad \StreamNameUp,\quad \RegSeqRatUp,\quad \RealUp.
$$
The same theorem keeps $\RealCompletionExactBoundaryUp$ as the boundary
witness on the same route rather than as a fifth roadmap object or an
aggregate status row.

Combining the two reads gives the public Phase Real and Completion exit
criterion from the Real chapter itself. The cover induction supplies the
finite Real terminal surface, the route exhaustion and minimality theorems
pin that surface to the L10 route, and the phase-exit projection consumes the
four local scoped/audit records while preserving the same-route exact-boundary
witness. No collapsed completion object, ambient completeness principle,
selected limit, quotient stream, countable choice, aggregate formalstatus
row, or bridge token is introduced by the composition.
\end{proof}

\begin{theorem}[Real L10 route-cover phase-exit converse]
\label{thm:real-l10-route-cover-phase-exit-converse}
Let $\mathcal{U}$ be a finite Real-facing replay-cover whose accepted public
read supplies the current Phase Real and Completion exit. Then the public exit
read pulls back exactly to the four object faces
$$
\DyadicRatCoreUp,\quad \StreamNameUp,\quad \RegSeqRatUp,\quad \RealUp.
$$
The pullback is witnessed by the displayed Real L10 support rows
$$
\BHist,\quad \hsame,\quad \Cont,\quad \Pkg,\quad \NameCert
$$
carried on the same route. Each of the four faces must carry the local
scoped-closure and audit-clean records required by
\autoref{thm:phase-real-completion-exit-four-object-threshold-projection};
the same-route $\RealCompletionExactBoundaryUp$ row remains a boundary witness
and not a fifth roadmap face.

Equivalently, no cover member, overlap, or public read of $\mathcal{U}$ can
serve the Phase Real and Completion exit unless its visible data are already
the four listed faces with those displayed support rows and those local
scoped/audit records. Quotient streams, selected limits, countable choice,
metric completeness, aggregate formalstatus rows, bridge tokens, and ambient
$\RealUp$ completeness are not admissible substitutes for any of the four
faces or their local records.
\end{theorem}
\begin{proof}
Start from the accepted public exit read and restrict it to the finite
Real-facing replay-cover. Apply
\autoref{thm:real-l10-route-cover-phase-exit-exactness} to identify the
available public read with the Real L10 terminal support surface whenever the
local threshold records are present.

It remains to show that no smaller or off-route surface can supply those
records. Apply \autoref{thm:real-l10-terminal-support-cover-induction} to
localize every cover contribution to a terminal support row, then apply
\autoref{thm:real-l10-terminal-four-face-support-minimality-pullback}. The
minimality pullback says that deleting any one of the
$\DyadicRatCoreUp$, $\StreamNameUp$, $\RegSeqRatUp$, or local $\RealUp$ faces,
or deleting one of the displayed $\BHist$, $\hsame$, $\Cont$, $\Pkg$, or
$\NameCert$ rows, loses the accepted terminal verdict.

The route itself is fixed by
\autoref{thm:real-l10-route-sibling-exhaustion}. Thus any attempted witness
outside the L10 route is pushed back to the same four-face support surface, or
else is rejected as an off-route consumer packet.

Finally apply
\autoref{thm:phase-real-completion-exit-four-object-threshold-projection}.
The phase-exit projection requires the four roadmap faces to be read through
their own local scoped-closure and audit-clean records, while keeping
$\RealCompletionExactBoundaryUp$ as a same-route boundary witness. The same
projection rules out quotient streams, selected limits, countable choice,
metric completeness, aggregate formalstatus rows, bridge tokens, and ambient
$\RealUp$ completeness as substitutes. Hence the public exit read pulls back
exactly to the four displayed faces with the displayed support rows.
\end{proof}

\begin{theorem}[Real L10 four-face public pullback criterion]
\label{thm:real-l10-four-face-public-pullback-criterion}
Let $\mathcal{P}$ be a finite consumer packet presented on the Real L10 route
at the public Phase Real and Completion boundary. Then $\mathcal{P}$ is
exactly the four-face pullback for that boundary if and only if its accepted
public data agree componentwise with the following four object faces:
\begin{enumerate}
\item the $\DyadicRatCoreUp$ endpoint, radius, and tolerance ledger;
\item the $\StreamNameUp$ finite-window rows;
\item the $\RegSeqRatUp$ finite-request, readback, and classifier rows;
\item the local $\RealUp$ terminal seal rows.
\end{enumerate}
The agreement is witnessed using only the displayed $\BHist$, $\hsame$,
$\Cont$, $\Pkg$, and $\NameCert$ support carried by the Real L10 route.
Equivalently, two such finite consumer packets determine the same public
Phase Real and Completion read precisely when they have the same four listed
faces and the same displayed support rows. The criterion does not use
quotient stream equality, selected host limits, countable-choice selectors,
ambient completeness, metric-completeness rows, aggregate formalstatus rows,
or bridge tokens.
\end{theorem}
\begin{proof}
For the forward direction, assume $\mathcal{P}$ is the public four-face
pullback. Apply
\autoref{thm:real-l10-route-cover-phase-exit-exactness} to read the public
boundary through $\DyadicRatCoreUp$, $\StreamNameUp$, $\RegSeqRatUp$,
$\RealUp$, and the same-route $\RealCompletionExactBoundaryUp$ witness.
Then apply \autoref{thm:real-l10-route-sibling-exhaustion} to remove any
off-route source, and apply
\autoref{thm:real-l10-terminal-four-face-support-minimality-pullback} to see
that the displayed $\BHist$, $\hsame$, $\Cont$, $\Pkg$, and $\NameCert$ rows
are the whole terminal support.

For the reverse direction, assume the four object faces and the displayed
support rows agree componentwise. The
\autoref{thm:streamname-regseqrat-realup-finite-window-pullback-criterion}
transports the StreamName finite-window rows through the RegSeqRat and local
Real readback surface, while
\autoref{thm:regseqrat-seal-route-visible-source-normal-form} normalizes the
finite-request, readback, classifier, and seal rows to the visible source
format. Finally
\autoref{thm:realup-open-phase-four-face-pullback} identifies the resulting
Real-facing packet with the four-face pullback at the open Phase Real
boundary.

Combining the two directions gives the stated if and only if. The exclusion
clauses are inherited from route exhaustion, terminal support minimality, the
route-cover exactness theorem, and the three finite-window and seal-route
pullback criteria just used. Each of those results keeps the construction on
finite NameCert support, so quotient stream equality, selected host limits,
countable choice, ambient or metric completeness, aggregate formalstatus
rows, and bridge tokens never enter the witness.
\end{proof}

\begin{theorem}[Real L10 terminal cover support normal form]
\label{thm:real-l10-terminal-cover-support-normal-form}
Let $\mathcal{U}$ be a finite family of accepted Real-facing terminal
presentations on the current Phase Real and Completion exit route. Then
$\mathcal{U}$ satisfies the terminal cover support normal form if and only if
each member of $\mathcal{U}$, and each overlap between two members of
$\mathcal{U}$, factors through exactly the four L10 faces
$\DyadicRatCoreUp$, $\StreamNameUp$, $\RegSeqRatUp$, and $\RealUp$, together
with the displayed $\BHist$, $\hsame$, $\Cont$, $\Pkg$, and $\NameCert$ rows
carried by those faces. No quotient stream, host real equality, selected
limit, choice selector, metric-completeness row, ambient completeness
coordinate, aggregate formalstatus row, or bridge row is part of this normal
form.
\end{theorem}
\begin{proof}
For the forward direction, assume that $\mathcal{U}$ satisfies the terminal
cover support normal form. Apply
\autoref{thm:real-l10-terminal-support-cover-induction} to reduce each member
and overlap of the finite cover to terminal support on the Real L10 route.
Then apply
\autoref{thm:streamname-finite-cover-terminal-support-minimality} and
\autoref{thm:regseqrat-four-face-finite-cover-normal-form}. These two
finite-cover results identify the StreamName and RegSeqRat contributions
with the same four displayed object faces and the same displayed support
rows.

It remains to account for the public Phase Real and Completion boundary.
\autoref{thm:real-l10-four-face-public-pullback-criterion} transports the
terminal cover data through the public pullback read. Since the pullback
criterion excludes off-surface witnesses, every accepted member and every
accepted overlap factors only through $\DyadicRatCoreUp$, $\StreamNameUp$,
$\RegSeqRatUp$, $\RealUp$, and the displayed $\BHist$, $\hsame$, $\Cont$,
$\Pkg$, and $\NameCert$ rows.

For the reverse direction, assume that every member and overlap of
$\mathcal{U}$ factors through exactly the four displayed faces with the
displayed support rows. Apply
\autoref{thm:real-l10-route-cover-phase-exit-exactness} to package this
four-face data as Real L10 route cover data at the Phase Real and Completion
exit. Then use
\autoref{thm:real-l10-route-cover-phase-exit-converse} to recover accepted
Real-facing terminal presentations from that route cover packet.

Finally apply
\autoref{thm:real-l10-four-face-public-pullback-criterion} once more to
identify the rebuilt terminal presentations with the public four-face
pullback. The excluded rows cannot be added during this reconstruction:
terminal support cover induction accounts only for finite support rows, the
StreamName and RegSeqRat normal forms account only for finite cover faces, and
the public pullback criterion keeps quotient streams, host real equality,
selected limits, choice selectors, metric or ambient completeness data,
aggregate formalstatus rows, and bridge rows outside the accepted normal
form.
\end{proof}

\begin{theorem}[Real L10 terminal route support bicriterion]
\label{thm:real-l10-terminal-route-support-bicriterion}
Let $\mathcal{U}$ be a finite Real-facing terminal cover at the current Phase
Real and Completion public route. Then $\mathcal{U}$ presents that public
route if and only if every member of $\mathcal{U}$, and every overlap between
two members of $\mathcal{U}$, factors through exactly
$\DyadicRatCoreUp$, $\StreamNameUp$, $\RegSeqRatUp$, local $\RealUp$, and the
displayed $\BHist$, $\hsame$, $\Cont$, $\Pkg$, and $\NameCert$ rows carried by
those four faces. Equivalently, two such terminal covers have the same public
route read precisely when their $\DyadicRatCoreUp$, $\StreamNameUp$,
$\RegSeqRatUp$, and local $\RealUp$ faces agree and their displayed support
rows agree.

No quotient stream, selected limit, countable-choice selector,
metric-completeness datum, aggregate formalstatus row, bridge row, or ambient
$\RealUp$ completeness witness is part of either side of the criterion.
\end{theorem}
\begin{proof}
For the forward direction, assume that $\mathcal{U}$ presents the public
Phase Real and Completion route. Apply
\autoref{thm:real-l10-terminal-cover-support-normal-form} to put the finite
terminal cover in support normal form. Then apply
\autoref{thm:real-l10-four-face-public-pullback-criterion} to identify the
public route read with equality of the four displayed object faces and the
displayed $\BHist$, $\hsame$, $\Cont$, $\Pkg$, and $\NameCert$ rows.

For the reverse direction, assume that every member and overlap factors
through exactly the four listed faces and the displayed support rows. Apply
\autoref{thm:real-l10-route-cover-phase-exit-exactness} to package these
finite rows as Real L10 route-cover data at the Phase Real and Completion
exit. Then apply
\autoref{thm:real-l10-route-cover-phase-exit-converse} to recover the
accepted terminal presentations on the same public route.

The equivalence of public route reads is the same argument with two covers in
parallel: the normal-form theorem reduces each cover to the four faces and
support rows, and the four-face public pullback criterion compares exactly
those data. The exclusions are preserved because each cited theorem keeps the
route finite and terminal: route-cover exactness and converse introduce only
accepted Real L10 rows, the normal form admits only the four displayed faces,
and the public pullback criterion refuses quotient streams, selected limits,
countable choice, metric or ambient completeness data, aggregate formalstatus
rows, and bridge rows.
\end{proof}

\begin{theorem}[Real L10 terminal route support obstruction]
\label{thm:real-l10-terminal-route-support-obstruction}
Let $\mathcal{U}$ be a finite Real-facing terminal cover at the current Phase
Real and Completion public route. If $\mathcal{U}$ does not present that public
route, then some member of $\mathcal{U}$, or some overlap between two members
of $\mathcal{U}$, fails to factor through one visible L10 face
$$
\DyadicRatCoreUp,\quad \StreamNameUp,\quad \RegSeqRatUp,\quad \RealUp
$$
or through one displayed Real L10 support row among
$$
\BHist,\quad \hsame,\quad \Cont,\quad \Pkg,\quad \NameCert.
$$
The obstruction is local to that missing face or support row. Quotient
streams, selected limits, countable-choice selectors, metric or ambient
completeness data, aggregate formalstatus rows, and bridge rows cannot repair
the failed public route presentation.
\end{theorem}
\begin{proof}
Argue contrapositively from
\autoref{thm:real-l10-terminal-route-support-bicriterion}. If every member
and every overlap of $\mathcal{U}$ factors through the four displayed faces
and the five displayed support rows, then the bicriterion presents the current
Phase Real and Completion public route.

Therefore a finite cover that does not present the public route must fail one
of the finite factorization requirements consumed by the bicriterion. Apply
\autoref{thm:real-l10-terminal-cover-support-normal-form} to put those
requirements in terminal-cover normal form. This identifies the possible
failure sites with the object rows
$\DyadicRatCoreUp$, $\StreamNameUp$, $\RegSeqRatUp$, and local $\RealUp$,
together with the displayed $\BHist$, $\hsame$, $\Cont$, $\Pkg$, and
$\NameCert$ rows.

It remains to rule out an off-surface rescue. Apply
\autoref{thm:real-l10-four-face-public-pullback-criterion} to the public
boundary read: the public route is determined exactly by the four visible
faces and the displayed support rows. Then apply
\autoref{thm:real-l10-route-cover-phase-exit-converse} to pull any accepted
public exit read back to the same Real L10 surface.

The two localization steps leave no constructor for quotient streams,
selected limits, countable-choice selectors, metric or ambient completeness
data, aggregate formalstatus rows, or bridge rows. Hence any failed finite
Real-facing terminal cover is obstructed at one visible face or at one
displayed support row, and no off-route datum repairs that failure.
\end{proof}
